\chapter{Results}
\label{chap:results}

This chapter reports the empirical findings of the baseline process-discovery and simulation-discovery pipeline. The aim is to assess whether the discovered process model is structurally valid, whether the simulation parameters can be inferred from historical data, and whether the resulting model provides a defensible basis for operational scenario analysis. The presentation follows the experimental pipeline of the thesis: process discovery, arrival discovery, resource discovery, duration discovery, what-if scenario analysis, and a discussion of the scientific interpretation of the results.

\section{Process Discovery Results}
\label{sec:process_discovery_results}

The first stage of the analysis assessed the quality of the discovered process model extracted from the cleaned event log. This is a necessary prerequisite for subsequent simulation modeling: without a structurally valid process model, parameter discovery would not be interpretable or reliable. Table~\ref{tab:process_discovery_quality} summarizes the core quality metrics.

\begin{table}[h!]
\centering
\caption{Process Discovery Quality Metrics}
\label{tab:process_discovery_quality}
\begin{tabular}{lr}
\toprule
Metric & Value \\
\midrule
Number of variants & 134 \\
Top-3 variant coverage & 67.48\% \\
Fitness & 1.000 \\
Precision & 0.903 \\
\bottomrule
\end{tabular}
\end{table}

The discovered model is highly consistent with the observed process behavior. A fitness value of 1.000 indicates that the Petri net can replay the entire event log without rejecting any recorded trace. In a process-mining context, fitness is the minimum requirement for a valid process model. If fitness were substantially below one, the model would fail to describe the observed operational behavior and would not provide a reliable basis for simulation generation.

The precision value of 0.903 is equally important. It indicates that the model is not merely permissive or over-generalized. Instead, it preserves a sufficiently constrained structure to distinguish realistic behavior from spurious or overly broad execution paths. This is particularly relevant in the CTB setting, where yard operations are constrained by physical resources, workload pressure, and routing decisions. The combination of perfect fitness and strong precision is therefore a strong indication that the process model is both complete and operationally meaningful.

The presence of 134 variants does not imply a structurally chaotic process. Rather, it indicates that the terminal operation is executed through many alternative but recurrent combinations of activities under different contextual conditions. The top-3 variant coverage of 67.48\% shows that a relatively small number of dominant behavioral patterns account for most of the process executions, while the remaining variants correspond to less frequent operational deviations. This is a common characteristic of logistics processes: many concrete executions are possible, but only a limited number of routings dominate the observed behavior.

Table~\ref{tab:petri_net_characteristics} reports the structural size of the discovered Petri net. The model contains 42 places, 59 transitions, and 128 arcs. These values suggest a process model that is rich enough to represent the operational flow of the terminal but still compact enough to remain interpretable and parameterizable. This balance is essential for simulation discovery: the model must be expressive enough to encode the system structure without becoming too large or opaque for valid parameter estimation.

\begin{table}[h!]
\centering
\caption{Petri Net Characteristics}
\label{tab:petri_net_characteristics}
\begin{tabular}{lr}
\toprule
Characteristic & Value \\
\midrule
Places & 42 \\
Transitions & 59 \\
Arcs & 128 \\
\bottomrule
\end{tabular}
\end{table}

These findings establish a sound foundation for the next stage of the analysis. The process model is not only structurally valid, but also operationally meaningful and well suited for parameter discovery. This motivates the simulation-discovery stage, which seeks to estimate the dynamic behavior governing timing, routing, and resource usage.

\section{Simulation Discovery Results}
\label{sec:simulation_discovery_results}

The simulation-discovery stage focused on learning the parameters that govern the dynamic execution of the process: arrivals, resource capacities, routing behavior, and contextual decision rules. Following the methodological framework of simulation discovery in process mining and data-aware simulation \cite{PMRozinat2009,Camargo2020,lopez2024,Vinci2026prosit,mannhardt2023stochasticProcessesModelling}, the project compared traditional pooled models against contextual models able to condition behavior on process-state variables.

\subsection{Arrival Discovery Results}
\label{sec:arrival_discovery_results}

Arrival behavior was the clearest empirical case for temporal and data-aware simulation. A pooled arrival model was compared against a time-aware decision tree conditioned on temporal process context. Table~\ref{tab:arrival_model_benchmark} summarizes the results.

\begin{table}[h!]
\centering
\caption{Arrival Model Benchmark}
\label{tab:arrival_model_benchmark}
\begin{tabular}{lrrr}
\toprule
Model & MAE & RMSE & $R^2$ \\
\midrule
Pooled arrival model & 29.23 & 35.81 & -0.007 \\
Time-aware decision tree & 11.66 & 19.56 & 0.700 \\
\bottomrule
\end{tabular}
\end{table}

The pooled arrival model performs poorly. Its RMSE of 35.81 and $R^2$ value of -0.007 indicate that a single global arrival rate provides essentially no explanatory power for the arrival process. By contrast, the time-aware decision tree reduces the MAE to 11.66 and the RMSE to 19.56, while achieving an $R^2$ of 0.700. This is a substantial improvement and provides clear evidence that arrival intensity is strongly conditioned by temporal context.

The day-vs-night statistics confirm the same result. Table~\ref{tab:arrival_day_night} compares the arrival behavior across the full sample, daytime, and nighttime periods.

\begin{table}[h!]
\centering
\caption{Day vs Night Arrival Statistics}
\label{tab:arrival_day_night}
\begin{tabular}{lrrr}
\toprule
Period & Mean interarrival (min) & Mean arrivals/hour & Cases \\
\midrule
All & 0.9539 & 62.90 & 89458 \\
Day & 0.5306 & 113.08 & 79309 \\
Night & 4.2618 & 14.08 & 10149 \\
\bottomrule
\end{tabular}
\end{table}

The difference between day and night operations is substantial: the mean arrival rate rises from 14.08 arrivals/hour at night to 113.08 arrivals/hour during the day, corresponding to an approximately eightfold difference. This is far beyond what a pooled arrival model can realistically represent. The empirical evidence therefore indicates that the arrival process is non-stationary and strongly time dependent. This is the strongest evidence in favor of data-aware simulation found in the project and provides a direct justification for modeling arrival behavior through temporal context.

\subsection{Resource Discovery Results}
\label{sec:resource_discovery_results}

The resource-discovery stage inferred effective capacities from overlap behavior observed in the historical logs. Because the event log contains pooled resource labels rather than explicit worker or equipment calendars, the resulting capacities should be interpreted as effective resource capacities rather than exact physical staffing counts. This distinction is essential for a scientifically sound interpretation of the results.

Table~\ref{tab:resource_capacities} reports the discovered capacities for the main resource pools.

\begin{table}[h!]
\centering
\caption{Discovered Resource Capacities}
\label{tab:resource_capacities}
\begin{tabular}{lrrrr}
\toprule
Resource & Median & P90 & Max & Capacity \\
\midrule
GateIn & 0 & 0 & 10 & 1 \\
GateOut & 0 & 0 & 9 & 1 \\
HO2 & 1 & 8 & 15 & 3 \\
LL & 1 & 6 & 25 & 3 \\
OTHER & 4 & 27 & 63 & 3 \\
\bottomrule
\end{tabular}
\end{table}

The overlap analysis shows that the larger pools are used concurrently, while gate-related resources remain largely single-threaded. This is operationally consistent with the terminal context: large handling areas support parallel execution, whereas gate operations remain tightly constrained by entry and exit procedures. For example, the OTHER pool shows a median concurrency of 4 and a maximum concurrency of 63, which demonstrates that the process clearly involves simultaneous handling within pooled resource categories.

However, these capacities are approximations. The historical event log does not include individualized resource calendars, worker schedules, or equipment assignments. Therefore, the simulation cannot recover the true physical resource plan from the log alone. Instead, it reconstructs the effective operational capacity that is observable in the data. This is a relevant limitation for the modeling exercise: the discovered resource model is valid for explaining observed overlap and operational congestion, but it cannot fully represent the exact physical resource allocation policy of the terminal.

\subsection{Data-Aware Duration Results}
\label{sec:data_aware_duration_results}

The duration-discovery step evaluated whether contextual process variables improve the predictive quality of activity-duration models with respect to a pooled baseline. The central finding is that the effect is selective rather than uniform. Table~\ref{tab:data_aware_durations} summarizes the activities for which contextual modeling showed meaningful explanatory value.

\begin{table}[h!]
\centering
\caption{Contextual Duration Model Performance by Activity}
\label{tab:data_aware_durations}
\begin{tabular}{lrr}
\toprule
Activity & $R^2$ & Interpretation \\
\midrule
LL\_delivery & 0.257 & Context helps meaningfully \\
LL\_mixed & 0.190 & Context helps moderately \\
HO2\_delivery & 0.133 & Context has some explanatory power \\
HO2\_receive & 0.137 & Context has some explanatory power \\
RMG\_delivery & 0.017 & Weak contextual effect \\
RMG\_receive & 0.005 & Essentially no contextual gain \\
\bottomrule
\end{tabular}
\end{table}

This table is important because it makes the actual story visible: context-aware simulation helps some activities much more than others. The strongest gains are observed for LL\_delivery and LL\_mixed, where the contextual model explains roughly 25.7\% and 19.0\% of the variance, respectively. Activities such as HO2\_delivery and HO2\_receive also show some explanatory value, but the effect is much weaker. For RMG-related activities, the contextual contribution is minimal: $R^2$ values are close to zero. This indicates that the behavior of certain handling stages is not sufficiently explained by the available contextual attributes alone.

This is a key methodological result. It shows that data-aware simulation is not automatically beneficial for all simulation components. Rather, its usefulness depends on whether the selected contextual features capture the true sources of variation for a given activity. In some cases, contextual variables such as target utilization, demand, and target rank explain a meaningful share of the variance; in others, the process remains dominated by unobserved operational rules, physical constraints, or local execution conditions. The thesis therefore treats data-aware simulation as a selective modeling improvement rather than as a universal replacement for distribution-based simulation.

\subsection{Feature Importance Analysis}
\label{sec:feature_importance_results}

To understand which contextual variables actually informed the tree-based models, the feature-importance analysis was examined for each activity. This analysis is important because it answers a central methodological question: which contextual variables are genuinely relevant for explaining process behavior, and which ones remain largely irrelevant in the discovered models? Table~\ref{tab:feature_importance_summary} summarizes the dominant feature patterns across the activity-specific trees.

\begin{table}[h!]
\centering
\caption{Dominant Features in Data-Aware Duration Models}
\label{tab:feature_importance_summary}
\begin{tabular}{ll}
\toprule
Activity & Dominant feature(s) \\
\midrule
LL\_delivery & hour, target utilization, target demand \\
LL\_mixed & hour, target utilization, target demand \\
HO2\_delivery & hour, target demand, weekday \\
HO2\_receive & hour, target utilization, weekday \\
RMG\_delivery & n\_deliveries, target utilization, visit complexity \\
RMG\_receive & target demand, visit complexity, hour \\
\bottomrule
\end{tabular}
\end{table}

The feature-importance analysis shows that the most relevant contextual variables are not uniform across the process. Time-related variables, especially hour, dominate several yard-handling activities, particularly LL and HO2 operations. For example, in LL\_delivery and LL\_mixed, hour contributes approximately 0.81 and 0.84 of the total importance, respectively, while target utilization and demand remain secondary but still relevant. Similarly, HO2\_delivery and HO2\_receive assign substantial importance to hour and to operational demand or utilization features.

This is consistent with the main arrival-discovery result. The process is strongly time dependent, and time of day appears to be one of the main drivers of operational variation in the yard. In contrast, target rank is largely irrelevant in most models, and some variables such as is\_weekend are consistently zero across the relevant activity trees. This suggests that the selected contextual features are informative only when they reflect operational pressure or time-dependent yard conditions, rather than static hierarchy or rank-based information.

The feature-importance analysis therefore closes an important methodological gap. It shows that the data-aware models are not arbitrary black-box predictors, but rather rely on a small set of operationally meaningful variables. At the same time, it confirms that the explanatory value of context depends on the activity and the process state. This is precisely why the thesis emphasizes selective usefulness: data-aware simulation is powerful when the process state contains relevant operational information, but it does not automatically improve every model component in the simulation bundle.

\section{Scenario Analysis Results}
\label{sec:scenario_analysis_results}

The final empirical analysis examined whether the discovered simulation model reacts plausibly to realistic operator interventions. Two what-if experiments were evaluated: a closure of block T22 and a second experiment combining T22 closure with a reduced RMG resource pool. In both cases, the objective was to test whether the white-box simulation model generated plausible operational adaptations rather than purely arbitrary outputs.

\begin{table}[h!]
\centering
\caption{Scenario Analysis: Baseline vs Resource Adaptation Experiments}
\label{tab:scenario_analysis}
\begin{tabular}{lrrr}
\toprule
Scenario & Mean turnaround (min) & Mean RMG service (min) & Mean RMG waiting (min) \\
\midrule
Baseline & 30.90 & 11.24 & 9.00 \\
T22 closed & 29.78 & 11.30 & 8.54 \\
T22 closed + reduced RMG pool & 29.99 & 10.85 & 8.32 \\
\bottomrule
\end{tabular}
\end{table}

The scenario results show that the model does not collapse under resource disruption. Instead, it adapts by redistributing workload across the remaining feasible blocks and resources. In the T22 closure scenario, the mean turnaround time decreases from 30.90 to 29.78 minutes and the mean RMG waiting time decreases from 9.00 to 8.54 minutes. This indicates that workload is rebalanced rather than causing a direct system-wide deterioration. Similarly, when the RMG resource pool is additionally reduced, the mean turnaround time remains close to the baseline at 29.99 minutes, while the mean RMG service and waiting times remain stable or slightly improved.

This behavior is important because it demonstrates the white-box nature of the discovered simulation. The model does not simply output a black-box forecast; it exposes how the process adapts to structural constraints through routing and resource reallocation. For operational decision support, this is highly valuable: it allows planners to reason about how a resource change affects traffic flow, congestion, and service times rather than only observing an aggregated KPI.

\section{Discussion}
\label{sec:results_discussion}

The results support a nuanced scientific conclusion. The process model is structurally excellent: it shows perfect replay fitness and acceptable precision, indicating that the historical truck process is well represented by the discovered Petri net. The arrival model shows a strong contextual dependence, which is the clearest empirical evidence in favor of data-aware simulation discovered in this project. Resource discovery is possible but constrained by the fact that the event log captures pooled resources rather than individualized worker or equipment schedules. Duration discovery is selective: contextual models help some activities substantially more than others, and the benefit is highly dependent on the activity and the available process variables. Finally, the what-if experiments show plausible adaptation behavior under resource changes, which makes the model useful for decision support in operational planning.

Taken together, these findings support the central claim of the thesis: data-aware simulation is valuable in this CTB setting, but only selectively and within the methodological boundaries of the available data. The value of data-aware modeling is strongest where the process exhibits clear temporal or state dependency, and weakest where operational variability is dominated by hidden constraints and pooled-resource dynamics. This is not a weakness of the method; it is an empirical characteristic of the real logistics process being modeled.

\section{Summary of Findings}
\label{sec:results_summary}

The empirical results show that the discovered process model is structurally valid and that the simulation-discovery pipeline is operationally meaningful. The arrival process is strongly time-dependent, and the contextual arrival model clearly outperforms the pooled baseline. Resource capacities can be inferred from overlap behavior but remain effective pooled capacities rather than exact physical resource plans. Data-aware duration models are useful for some activities, yet their gain varies substantially across the process. The white-box simulation responds plausibly to resource disruptions, confirming that the model is suitable for scenario analysis and decision support under realistic operational changes.

Overall, the results indicate that a data-aware process simulation approach is promising for CTB truck operations, but its benefits are conditional on the activity type, the modeling objective, and the degree to which the available historical data represent the underlying physical system.
